\subsection{Shared search interface and coupled parameter profile}
\label{sec:room-approach}

The room adapter reuses the five searches on waiting-case-to-room allocations
and evaluates each through the decoder and Equation~\ref{eq:room-cost}.
Reassignment and swap moves operate on this common space. At opening, a
greedy baseline inserts cases by increasing prediction into the earliest
available room. At replanning, the incumbent retains room choices from the
accepted plan and greedily inserts previously unassigned cases before decoding
against current availability. Local searches start from this incumbent;
GA includes it in the population and ACO retains it as the initial best.
Keeping the best evaluated candidate protects predicted request cost, while
execution can still deteriorate under duration error.

\begin{table}[htbp]
\centering\small
\caption{Coupled room-study settings. Iteration ceilings are raised so the
common fitness-call budget controls nontrivial searches. These fixed settings
are not the outcome of a tuning experiment.}
\label{tab:room-settings}
\begin{tabularx}{\linewidth}{@{}lY@{}}
\toprule
Method & Parameters \\
\midrule
SA & $T_0=1$, cooling $0.95$. \\
Tabu & Memory 20, 15 sampled neighbours. \\
GA & Population 40, crossover $0.8$, mutation $0.08$, elitism. \\
Hybrid & Memory 20, 15 neighbours, $T_0=1$, cooling $0.97$. \\
ACO & 15 ants, $\alpha=1$, $\beta=2$, evaporation $0.3$, deposit scale 1. \\
\bottomrule
\end{tabularx}
\end{table}

The room profile adapts both annealing temperatures to its normalised cost;
it must not be confused with the aggregate hybrid's $T_0=20$ default. ACO
constructs room assignments with probabilities proportional to
$\tau_{ir}^{\alpha}\eta_r^{\beta}$, where
$\eta_r=(1+L_r/c_r)^{-1}$ uses accumulated predicted load and remaining
capacity. This construction heuristic excludes turnover; complete candidates
are assessed by the common decoder. Pheromone evaporates by a factor $0.7$
per iteration, and the best ant's assignments receive $1/(1+C_{\mathrm{room}})$.
Pheromone is reinitialised for each request.

\subsection{Mesa execution and proposal acceptance}
\label{sec:coupling}

One Mesa agent represents a patient episode with waiting, running and
completed states; rooms are capacity resources. Coordination is centralised:
agents do not negotiate schedules or run their own optimisers. The execution
loop is as follows:
\begin{enumerate}
\item Build a request from the observable state, predictions and fixed
      running/completed assignments; omit hidden realised durations.
\item Obtain an allocation from the scheduler, decode it and check coverage,
      timing, turnover, known closures and preservation of fixed work.
\item Accept only a feasible proposal based on the current state version.
      A new state invalidates stale proposals; failed requests retain the
      accepted plan, or an all-unassigned plan if initial planning fails.
\item Advance one simulated minute: complete operations, release turnover,
      reveal closure/reopening, synchronise state and, under the reactive
      policy, request a new plan at those availability events.
\item Start eligible cases from the accepted schedule, subject to actual
      room occupancy, turnover, the start cutoff and closure guards.
\end{enumerate}
Static and reactive replays begin from the same saved initial schedule.
Duration overruns alone do not trigger replanning. A planned start remains
a lower bound until changed by an accepted replan. Execution duration is
unchanged by reassignment; ongoing work continues through a room closure.

Nontrivial searches run in separate processes with progress incumbents.
Budgets count initialisation and repeated fitness calls; cancellation/deadline
handling terminates and joins workers. A pre-search deadline may return a
validated incumbent with zero evaluations, labelled separately. Empty and
singleton requests use trivial/direct handling. Runtime includes startup and
communication; simulated time pauses during solving. This supports controlled
experiments, not a production response-time guarantee.
